\section{Studying strange hadron production in hard and soft processes}

In hadronic collisions, particle production mechanisms can be categorised into either hard or soft processes depending on the momentum transfer 
involved. Hard scattering processes involve high momentum transfer from an energetic parton resulting in a narrow, collimated spray of hadrons 
retaining a significant amount of the initial beam momentum, called a "jet". The charged hadron with the highest transverse momentum 
($p_{\mathrm{T}}$) in the event is called the jet-leading particle, or the "trigger" particle, which also acts as a proxy for the 
jet axis and defines the near-side jet region. Alternatively, soft processes involving particle production governed by bulk properties of the medium are 
characterised by low momentum transfer and can be studied in the transverse to leading or out-of-jet regions. Understanding the relative contribution of hard and soft 
processes to strangeness production in small collision systems is of particular interest and can be studied using techniques such as full-jet 
reconstruction and/or two-particle angular correlations. Recent ALICE results presented in these proceedings utilize the angular correlation method 
to distinguish strange hadrons ($\phi,\:\Lambda,\:\ks,\:\Xi^{\pm}$) produced in hard processes (e.g., jets) from soft processes (out-of-jet).

\subsection{Two-particle angular correlations}
Studying strange hadron production using the method of two-particle angular correlations \cite{Correlations:2013} with respect to jet axis in the 
($\Delta\varphi,\:\Delta\eta$) plane involves selecting the trigger particle and identifying strange hadrons of interest in the event. The analyses 
presented here impose kinematic selection criteria for the trigger particles: \pt$ > 3\, \mathrm{GeV}/c$ ($\ks,\,\Xi^{\pm}$) and 
$4 < $ \pt $< 8\, \mathrm{GeV}/c$ ($\phi,\,\Lambda$). From the same event, strange hadrons of interest with smaller transverse momenta (\pt) than the 
trigger particle are labeled as "associated" particles. The angular correlations can be constructed using the angular differences ($\Delta\varphi,\:\Delta\eta$) between the trigger 
and the associated particles. The angular distribution is divided into three regions: towards leading (near-side jet), 
transverse to leading (out-of-jet or underlying event), and full inclusive region, covering the entire angular correlation plane. For each region, 
per-trigger yield of associated particles is computed from the 2D correlation distribution and corrected for factors such as trigger and associated 
particles' reconstruction efficiencies, pair detector acceptance, and contamination from non-primary particles.